% ---
% id: meat-core/lem:dk
% kind: lem
% title: "Divided-difference derivative of a matrix function"
% summary: >-
%   At a Hermitian operator, a matrix function continuously differentiable
%   around its spectrum has as derivative its divided-difference kernel between
%   spectral projectors, self-adjoint under the trace pairing. At trace level
%   only the diagonal kernel survives.
% prominence: technical
% status: implemented
% uses: []
% mentions: [thm:gradient, lem:cone-concavity, lem:b-derivative]
% assumes: []
% refs: []
% provenance: newly-derived
% implements: certiqs_key.meat_core.objective.linalg:divided_difference
% lean:
% ---

Theorem~\ref{thm:gradient} differentiates $X\mapsto X^{p}$ inside the
support of a branch operator with the kernel below, and
Lemma~\ref{lem:b-derivative} uses its trace form.

\begin{lemma}\label{lem:dk}
Write $p\coloneqq1/\alpha\in(0,1)$ and $T$ for the functional
\eqref{eq:Tcone} of Lemma~\ref{lem:cone-concavity} on the positive
semidefinite cone of a branch space; $\log_2\mathrm{e}=1/\ln2$.

Let $X$ be Hermitian with spectral decomposition $X=\sum_i\lambda_iP_i$, and
let $f$ be continuously differentiable on an open interval
$I\supseteq\operatorname{spec}X$. Then $X'\mapsto f(X')$, defined on the
Hermitian operators with spectrum in $I$, is Fr\'echet differentiable at $X$
with
\begin{equation}\label{eq:dk}
  Df(X)[H]=\sum_{i,j}f^{[1]}(\lambda_i,\lambda_j)\,P_iHP_j,
  \qquad
  f^{[1]}(\lambda,\mu)\coloneqq
  \begin{cases}\dfrac{f(\lambda)-f(\mu)}{\lambda-\mu},&\lambda\ne\mu,\\[1ex]
  f'(\lambda),&\lambda=\mu,\end{cases}
\end{equation}
and for Hermitian $A$, $H$:
$\tr\!\big[A\,Df(X)[H]\big]=\tr\!\big[Df(X)[A]\,H\big]$. In particular
this applies to $f(x)=x^{p}$ at every $X\succ0$, with $I=(0,\infty)$; and,
at the trace level, for every Hermitian $Y$ and every $g$ continuously
differentiable on an open interval containing $\operatorname{spec}Y$, the
map $Y'\mapsto\tr g(Y')$ is Fr\'echet differentiable at $Y$ with
\begin{equation}\label{eq:trace-diff}
  \frac{\mathrm d}{\mathrm dt}\Big|_{t=0}\tr\,g(Y+tK)=\tr\big[g'(Y)K\big]
  \qquad\text{for every Hermitian }K.
\end{equation}
For $\alpha>1$ the function $g(y)\coloneqq\max\{y,0\}^{\alpha}$ is
continuously differentiable on $\mathbb{R}$, with
$g'(y)=\alpha\max\{y,0\}^{\alpha-1}$, and $\tr g(Y)=\tr[Y^{\alpha}]$ for
every $Y\succeq0$.
\end{lemma}

\begin{proof}
\begin{sloppypar}
\emph{Polynomials.} For a monomial $f(x)=x^{m}$,
$\frac{\mathrm d}{\mathrm dt}(X+tH)^{m}\big|_{0}
=\sum_{a+b=m-1}X^{a}HX^{b}$, and
$\bra{i}\sum_{a+b=m-1}X^aHX^b\ket{j}
=H_{ij}\sum_{a+b=m-1}\lambda_i^a\lambda_j^b
=H_{ij}\,f^{[1]}(\lambda_i,\lambda_j)$, which is \eqref{eq:dk}; linearity
extends it to polynomials. A polynomial map $X'\mapsto q(X')$ is a finite
sum of products of $X'$ with itself, hence Fr\'echet differentiable at every
$X$, and its Fr\'echet derivative is the directional derivative just
computed.

\emph{A Hilbert--Schmidt bound.} Let $h$ be continuously differentiable on a
compact interval $[c,C]$, and let $A=\sum_ia_iP_i$ and $B=\sum_jb_jP'_j$ be
Hermitian with spectra in $[c,C]$. Then
$P_i\big(h(A)-h(B)\big)P'_j=(h(a_i)-h(b_j))P_iP'_j$ and
$P_i(A-B)P'_j=(a_i-b_j)P_iP'_j$, so
$h(A)-h(B)=\sum_{i,j}h^{[1]}(a_i,b_j)\,P_i(A-B)P'_j$, the terms with
$a_i=b_j$ being $0$ on both sides. The maps $Z\mapsto P_iZP'_j$ are
orthogonal projections of the Hilbert--Schmidt space onto mutually
orthogonal subspaces that sum to it, and $|h^{[1]}|\le\|h'\|_{\infty}$ on
$[c,C]^2$ by the mean value theorem, so
$\|h(A)-h(B)\|_2\le\|h'\|_{\infty}\,\|A-B\|_2$; in the same way
$\|\sum_{i,j}h^{[1]}(\lambda_i,\lambda_j)P_iHP_j\|_2
\le\|h'\|_{\infty}\,\|H\|_2$.

\emph{General $f$.} Choose a compact $[c,C]\subset I$ and $\delta>0$ with
$\lambda_i\pm\delta\in[c,C]$ for every $i$. By Weyl's inequality every
eigenvalue of $X+H$ lies within $\|H\|$ of one of $X$, so
$\operatorname{spec}(X+H)\subseteq[c,C]$ whenever $\|H\|_2\le\delta$. Pick
polynomials $q_n\to f'$ uniformly on $[c,C]$ (Weierstrass) and let $Q_n$ be
primitives with $Q_n(c)=f(c)$, so that
$\varepsilon_n\coloneqq\sup_{[c,C]}|f'-Q_n'|\to0$. Write
$L(H)\coloneqq\sum_{i,j}f^{[1]}(\lambda_i,\lambda_j)P_iHP_j$ and $L_n$ for
the same expression with $Q_n$, which is $DQ_n(X)$ by the polynomial step.
For $\|H\|_2\le\delta$ the Hilbert--Schmidt bound, applied to $h=f-Q_n$,
whose divided differences are those of $f$ less those of $Q_n$, gives
\begin{equation*}
  \big\|\big(f(X+H)-f(X)-L(H)\big)
  -\big(Q_n(X+H)-Q_n(X)-L_n(H)\big)\big\|_2
  \;\le\;2\varepsilon_n\|H\|_2 .
\end{equation*}
Given $\eta>0$, fix $n$ with $2\varepsilon_n\le\eta/2$ and then
$\delta'\in(0,\delta]$ with
$\|Q_n(X+H)-Q_n(X)-L_n(H)\|_2\le\frac{\eta}{2}\|H\|_2$ for
$\|H\|_2\le\delta'$; then $\|f(X+H)-f(X)-L(H)\|_2\le\eta\|H\|_2$ for
$\|H\|_2\le\delta'$. The estimate holds uniformly in $H$ on that ball, so
$f$ is Fr\'echet differentiable at $X$ with $Df(X)=L$, which is
\eqref{eq:dk}.

\emph{Self-adjointness.} With $k_{ij}\coloneqq
f^{[1]}(\lambda_i,\lambda_j)=k_{ji}$ real,
$\tr[A\sum_{ij}k_{ij}P_iHP_j]=\sum_{ij}k_{ij}\tr[P_jAP_iH]
=\tr\big[\big(\sum_{ij}k_{ji}P_iAP_j\big)H\big]
=\tr[Df(X)[A]\,H]$.

\emph{Trace level.} The general step applies to $g$ at $Y=\sum_i\mu_iP_i$,
so $Y'\mapsto g(Y')$ is Fr\'echet differentiable at $Y$ with derivative
$\sum_{i,j}g^{[1]}(\mu_i,\mu_j)P_iKP_j$, and so is its composition with the
linear trace. Only the diagonal kernel survives the trace:
$\tr[P_iKP_j]=\delta_{ij}\tr[P_iK]$, so
$\frac{\mathrm d}{\mathrm dt}\tr g(Y+tK)\big|_0
=\sum_ig'(\mu_i)\tr[P_iK]=\tr[g'(Y)K]$. For $\alpha>1$ the function
$g(y)=\max\{y,0\}^{\alpha}$ has derivative $\alpha y^{\alpha-1}$ at $y>0$,
$0$ at $y<0$, and $0$ at $y=0$, since $|g(y)|/|y|\le|y|^{\alpha-1}\to0$;
that derivative is continuous because $\alpha-1>0$; and $g$ is $y^{\alpha}$
on $[0,\infty)$, so $\tr g(Y)=\tr[Y^{\alpha}]$ for $Y\succeq0$.
\end{sloppypar}
\end{proof}
